% paper.tex —— 由 source/main.tex 拍平：\input 已展开，注释、废稿（\iffalse）、参考文献已去掉。
% 读正文就读这个文件；行号和 outline.md 对得上。
% ==== 正文用到的自定义宏（3 个；看不懂的记号先查这里）====
\newcommand{\stddev}[1]{\textcolor{gray}{\scalebox{.8}{$\pm$#1}}}
\newcommand{\circlednum}[1]{\tikz[baseline=(char.base)]{\node[shape=circle,draw=teal,fill=teal!15,inner sep=1pt,text=teal!90!black] (char) {#1};}}
\renewcommand{\maketitle}{ \thispagestyle{firstpage} \begin{tcolorbox}[ breakable, colback=lightbluebg, colframe=blueframe, arc=3mm, boxrule=0.5pt, left=12pt, right=12pt, top=12pt, bottom=12pt, width=\textwidth, boxsep=5pt ] \noindent {\sffamily\Large\bfseries\color{darkblue}\@title\par} \vspace{0.6em} \noindent \parbox{\textwidth}{\@author}\par \vspace{1.2ex} {\color{dividercolor}\rule{\linewidth}{0.5pt}}\par \vspace{1.2ex} \ifx\@abstract\@empty\else \noindent\parbox{\textwidth}{ {\sffamily\bfseries Abstract:}\quad \@abstract }\par \vspace{1ex} \fi \ifx\@keywords\@empty\else \noindent\parbox{\textwidth}{ {\sffamily\bfseries Keywords:}\quad \@keywords }\par \vspace{1.5ex} \fi \ifx\@codelink\@empty \ifx\@datasetlink\@empty \ifx\@projectlink\@empty \else \vspace{-1.1em} {\color{dividercolor}\rule{\linewidth}{0.5pt}} \vspace{-0.6em} \centering \sffamily\small \href{\@projectlink}{\faGlobe\ Project Page} \par \fi \else \vspace{-1.1em} {\color{dividercolor}\rule{\linewidth}{0.5pt}} \vspace{-0.6em} \centering \sffamily\small \ifx\@projectlink\@empty\else \href{\@projectlink}{\faGlobe\ Project Page}\hspace{1.5em} \fi \href{\@datasetlink}{\faDatabase\ Dataset} \par \fi \else \vspace{-1.1em} {\color{dividercolor}\rule{\linewidth}{0.5pt}} \vspace{-0.6em} \centering \sffamily\small \href{\@codelink}{\faGithub\ Code} \ifx\@datasetlink\@empty\else \hspace{1.5em}\href{\@datasetlink}{\faDatabase\ Dataset} \fi \ifx\@projectlink\@empty\else \hspace{1.5em}\href{\@projectlink}{\faGlobe\ Project Page} \fi \par \fi \end{tcolorbox} }
% ==== 正文 ====
\maketitle

\section{Introduction}
Multi-agent systems, powered by large language models (LLMs), have emerged as a promising paradigm for complex medical decision-making~\cite{li2023camel, arora2023promise,zhu2025healthflow}. By simulating interactions such as expert debates~\cite{du2023improving} or role-playing multidisciplinary teams~\cite{tang2024medagents,kim2024mdagents}, these systems aim to enhance diagnostic accuracy and clinical reasoning. Preliminary results are compelling, with frameworks reporting significant performance gains on medical question-answering (QA) benchmarks over single-agent baselines~\cite{chen2024reconcile,wu2024autogen,moor2023foundation}.

However, this enthusiasm must be tempered by a critical gap in current evaluation practices. Existing studies predominantly assess these systems through final-answer accuracy, treating the intermediate collaborative process as an opaque ``black box''. This overlooks a core question central to clinical practice and public trust: \emph{is a diagnostic conclusion, whether correct or not, reached through a sound and verifiable reasoning pathway?} In high-stakes medical settings, particularly as these systems are considered for deployment in clinical or public-facing applications, the transparency and reliability of the decision-making process are as crucial as the final outcome~\cite{ribeiro2016should}. An answer that is incidentally correct but derived from flawed logic or suppressed dissent offers no reliable assurance of future performance and cannot earn clinical trust~\cite{sbaffi2017trust}.

Recent work has begun to scrutinize the assumed benefits of multi-agent collaboration, identifying inherent process vulnerabilities like consensus formation around incorrect beliefs and the suppression of valuable minority viewpoints in general-purpose systems~\cite{pan2025multiagent,choi2025debate,zhu2025medagentboard}. Some studies even suggest that simpler ensembling methods can account for most performance gains~\cite{choi2025debate,kaesberg2025voting}, questioning whether observed benefits stem from genuine synergistic reasoning~\cite{madaan2023selfrefine}. These concerns are amplified in medicine, where failures in collaborative reasoning can have severe consequences~\cite{huang2023large}.

To bridge this gap, we present a comprehensive empirical investigation into the reliability of multi-agent collaboration in medical AI. We systematically analyze 3,600 interaction logs across six state-of-the-art frameworks and six diverse medical datasets. Our analysis reveals several dominant failure patterns: (1) loss of key correct information, where critical details are omitted during synthesis; (2) suppression of valuable minority opinions, where majority bias silences correct dissenting views; (3) bypassing of evidence-based evaluation, as decisions default to voting instead of argument quality; (4) loss of collaborative diversity, where role assignments fail to elicit domain-specific expertise; (5) failure to prioritize high-risk clinical outcomes, resulting in misleading diagnoses; and (6) self-contradictory output due to a lack of cross-turn memory. These findings reveal deep-seated limitations in current collaboration designs. The contributions of this work are as follows:
\begin{itemize}
    \item We provide a large-scale empirical analysis of multi-agent collaboration dynamics in medical AI, revealing recurring patterns of how agents interact, converge, and fail.
    \item We develop a structured taxonomy of collaborative failures through rigorous qualitative analysis, validated with high inter-annotator agreement (Cohen's $\kappa=0.82$), providing a structured vocabulary for future diagnostics.
    \item We introduce a quantitative auditing framework to trace the propagation of information, quantify how framework designs are associated with specific failure patterns, and identify mechanisms by which collaboration breaks down.
\end{itemize}

By systematically diagnosing the failures of current multi-agent medical systems, this work lays the groundwork for developing the transparent and reliable collaborative AI that clinical practice demands.

\section{Related Work}

\paragraph{The rise and scrutiny of multi-agent systems.}
Multi-agent systems~\cite{kim2024health} leverage multiple LLM instances to solve complex tasks through structured interactions~\cite{tran2025mas_survey} such as debate~\cite{du2023improving,hong2024metagpt}, role specialization~\cite{tang2024medagents}, and consensus formation~\cite{chen2024reconcile}. In medicine, systems like MedAgents~\cite{tang2024medagents} and ColaCare~\cite{wang2025colacare} have demonstrated performance gains on benchmark tasks. However, a growing body of work challenges the assumption that collaboration inherently improves outcomes. Recent studies reveal that simple ensembling methods like majority voting often account for most of the performance gains attributed to complex debate mechanisms~\cite{choi2025debate, kaesberg2025voting, liang2023encouraging}. Our work builds on this critical perspective, moving from asking if these systems work to diagnosing why and how they fail, specifically within the high-stakes medical domain.

\paragraph{Failure analysis in AI systems.}
Understanding failure is critical for building robust AI~\cite{arrieta2020explainable}. Research has characterized single-agent LLM failures like hallucination and reasoning errors~\cite{qiu2025quantifying,huang2025survey}. More recently, attention has shifted to multi-agent systems. For instance, a general taxonomy of MAS failures has been developed through grounded theory~\cite{pan2025multiagent}, and methods for automated failure attribution have also been introduced~\cite{zhang2025which}. Similarly, an empirical study was conducted on failures in automated software issue solving~\cite{liu2025empirical, cemri2025multi}. While these studies provide a valuable starting point, they are not tailored to the unique challenges of medicine, such as interpreting multimodal data, handling clinical uncertainty, and prioritizing patient safety. Our work provides a failure taxonomy specifically for medical multi-agent collaboration, capturing domain-specific failure modes like misinterpretation of visual evidence and errors in clinical risk assessment~\cite{esteva2019guide}.

\section{Methodology}
Our research follows a multi-phase design to systematically investigate and quantify collaborative failure modes in medical multi-agent systems, as illustrated in \Cref{fig:overview_methodology}.

\begin{figure*}[!ht]
    \centering
    \includegraphics[width=0.95\linewidth]{figures/overview_methodology.pdf}
    \caption{Overview of the four-phase research methodology, encompassing (A) data generation, (B) qualitative analysis and taxonomy development, (C) quantitative auditing design, and (D) quantitative auditing and reporting.}
    \label{fig:overview_methodology}
\end{figure*}

\subsection{Data Generation and Instrumentation}
The foundation of our study is a large-scale dataset of multi-agent interactions. We execute six distinct multi-agent frameworks on six medical datasets, generating a total of 3,600 complete interaction logs. Each log serves as a detailed transcript of a collaborative session.

\paragraph{Instrumentation for process auditing.}
To enable transparent analysis, we instrument the source code of each framework to produce a structured audit trail for every case. This trail captures the entire decision-making process in a machine-readable format. For each agent's turn, we record: (1) the full prompt sent to the LLM, including system and user messages; (2) the verbatim response from the LLM, alongside any extracted structured data; and (3) contextual metadata, such as the agent's identity, role, and the current interaction round. This instrumentation provides the raw data for all subsequent qualitative and quantitative analyses.

\subsection{Qualitative Taxonomy Development}
To understand and categorize the observed failures, we construct a taxonomy using a rigorous qualitative analysis grounded in established research practices~\cite{strauss1990basics,adolph2011using}.

\paragraph{Pilot analysis and initial codebook development.}
The process begins with a pilot analysis of a diverse sample of 300 logs, drawing 50 from each of the six frameworks. Two annotators, both with expertise in AI and medical informatics, independently analyze these logs. Following open coding principles~\cite{strauss1990basics}, they identify and assign descriptive labels to recurring failure patterns without preconceived categories. Subsequently, the annotators convene to compare their findings, consolidate overlapping labels, and refine the definitions of each category. This collaborative process results in a preliminary codebook that serves as the structured foundation for the full-scale annotation phase.

\paragraph{Full-scale annotation and taxonomy refinement.}
Next, we expand the annotation to a larger sample of 3,600 logs. Using the preliminary codebook, the two annotators independently code each case, assigning one or more failure labels and noting the specific interaction round where the critical failure occurred. Disagreements are resolved through regular discussions, leading to iterative refinements of the codebook's definitions and hierarchy. This cycle of coding, comparison, and refinement continues until theoretical saturation is reached---that is, when analyzing additional cases yields no new failure categories. This process results in a stable, hierarchical taxonomy of collaborative failure modes.

\paragraph{Reliability and validation.}
To ensure the objectivity of our taxonomy, we conduct an inter-annotator agreement (IAA) study on a held-out set of 200 randomly sampled logs. The two annotators independently code this set using the final taxonomy. We calculate the Cohen's Kappa coefficient~\cite{cohen1960coefficient}, which yields a score of 0.82, indicating almost perfect agreement. This result confirms that our taxonomy is a reliable instrument for identifying and categorizing collaborative failures.

\subsection{Quantitative Auditing Design}
To provide a conceptual framework for our analysis, we define the core agent archetypes within the systems under study. Domain agents are content producers, responsible for initial analysis and evidence extraction from the case data. Meta agents act as process managers, orchestrating the collaboration by synthesizing opinions, mediating disagreements, and making final decisions. Finally, Audit agents serve as non-participatory observers that execute our quantitative auditing mechanisms. Building on the qualitative insights from our taxonomy, these four mechanisms are designed to measure specific, critical aspects of the collaborative process (\Cref{fig:audit_mechanism_illustration}) by transforming qualitative failure concepts into scalable, objective metrics~\cite{lundberg2017unified,zheng2023llm_as_a_judge}, generating an audit trail for each interaction.

\begin{figure*}[!ht]
    \centering
    \includegraphics[width=0.95\linewidth]{figures/audit_mechanism_illustration.pdf}
    \caption{An overview of the AuditTrail framework, comprising four mechanisms designed to quantify and record the multi-agent collaborative process.}
    \label{fig:audit_mechanism_illustration}
\end{figure*}

\paragraph{Tracking of key evidential units.}
This mechanism measures information loss by tracking the propagation of critical facts. First, domain agents extract evidential units from the case. Second, an auditor agent identifies a subset of key evidential units (KEUs) necessary for a correct diagnosis. Third, we track whether these KEUs are present in the reasoning of subsequent synthesis and decision steps. This process allows us to calculate a KEU retention rate.

\paragraph{Attribution of viewpoint shifts.}
To diagnose failures such as suppression of valuable minority opinions, we analyze the dynamics of opinion change. We categorize each case into one of four patterns: (1) Successful Minority Correction (M1), where an initially incorrect majority is overturned by a correct minority; (2) Negative Majority Assimilation (M2), where a correct minority succumbs to an incorrect majority; (3) Robust Majority Resilience (M3), where a correct majority withstands incorrect minority opinions; and (4) Minority-Induced Derailment (M4), where an incorrect minority misleads a correct majority. By instrumenting agents to declare whether their opinion shifts are based on evidence or consensus, we quantify the prevalence of these dynamics and pinpoint whether collaboration fosters correction or conformity.

\paragraph{Evaluation of collaboration quality.}
This mechanism addresses failures related to the quality of the reasoning process. An auditor agent scores interactions along three dimensions: (1) Evidence-based decision-making: The auditor assigns an overall argument quality score to each agent's argument. A failure is detected if the system's final answer aligns with the majority vote rather than the argument with the highest quality score. (2) Role effectiveness: The auditor assesses the specialized insight emergence of each domain agent's contribution, quantifying the degree to which role-playing elicits diverse, expert knowledge. (3) Clinical risk management: The auditor classifies the diagnostic urgency level implied by each agent's reasoning, allowing us to measure the clinical priority mismatch rate where the system opts for a low-urgency diagnosis when a safer, high-urgency alternative was proposed.

\paragraph{Tracking of conflict resolution.}
To measure the system's ability to handle internal disagreements, this mechanism measures the system's ability to handle internal disagreements. The process is as follows: (1) After the initial round, an auditor agent identifies substantive contradictions between agents' arguments, labeling them as Critical Conflict Points (CCPs). (2) These CCPs are injected into subsequent prompts to compel agents to address them. (3) We track each CCP, to determine if it was substantively addressed or ignored. The primary metric is the conflict resolution dropout rate—the percentage of CCPs that remain unaddressed, which quantifies the system's tendency to bypass internal logical conflicts.

\section{Experimental Setups}

\subsection{Datasets and Frameworks}

\paragraph{Medical datasets.}
Our study spans six datasets to ensure diversity in both medical tasks and data modalities. For textual question answering, we select MedQA~\cite{jin2021medqa}, a multiple-choice dataset based on medical licensing exams; PubMedQA~\cite{jin2019pubmedqa}, which requires yes/no answers to questions based on biomedical research abstracts; and MedXpertQA~\cite{zuo2025medxpertqa}, a dataset designed to test expert-level medical knowledge. For visual question answering (VQA), we use PathVQA~\cite{he2020pathvqa}, focusing on pathology images; VQA-RAD~\cite{lau2018VQA_RAD}, based on clinical radiology images; and SLAKE~\cite{liu2021slake}, a semantically labeled radiology dataset. In total, our analysis covers 3,600 medical cases, providing a comprehensive basis for our investigation.

\paragraph{Medical multi-agent frameworks.}
We select six representative multi-agent frameworks that implement diverse collaboration mechanisms, allowing for a broad analysis of interaction dynamics. The audited frameworks include two consensus-based systems, ColaCare~\cite{wang2025colacare} and ReConcile~\cite{chen2024reconcile}, which focus on iterative refinement and agreement. We also evaluate two role-playing frameworks, MDAgents~\cite{kim2024mdagents} and MedAgents~\cite{tang2024medagents}, where agents assume specialized medical personas. To cover other common paradigms, we include MAC~\cite{chen2025enhancing}, a hierarchical system employing a supervisor agent, and HealthcareAgent~\cite{ren2025healthcare}, a standard debate-based framework~\cite{du2023improving}. This selection covers the dominant architectural patterns in current multi-agent research.

\subsection{Implementation Details}

All experiments are conducted on a Mac Studio M3 Ultra with 512GB of RAM. We develop this work using Python 3.12. For all text-based QA tasks, generative agents responsible for analysis and synthesis are powered by DeepSeek-V3.2~\cite{liu2024deepseekv3,deepseek_v32}. For VQA tasks, agents that directly process multimodal input utilize Qwen3-VL~\cite{Qwen25VL,Qwen3VL} (\texttt{qwen3-vl-235b-a22b-instruct}), a state-of-the-art, open-source large vision-language model. In frameworks with heterogeneous roles, text-only meta-agents (e.g., synthesizers or deciders) also use DeepSeek-V3.2 on VQA tasks, operating on the textual outputs provided by the vision-capable agents. Critically, to maintain the highest level of analytical objectivity and minimize bias, all our diagnostic auditor agents, including those executing the four quantitative auditing mechanisms, are exclusively powered by DeepSeek-V3.2-Think~\cite{liu2024deepseekv3,guo2025deepseekr1}, a model selected for its strong reasoning capabilities.

\section{A Taxonomy of Collaborative Failure Modes}
Our qualitative analysis of 3,600 interaction logs culminates in a comprehensive taxonomy that reveals a critical disconnect between a task's outcome and its underlying collaborative process. A key finding is that even cases with correct final answers often suffer from deeply flawed collaboration. Many successes are not the outcome of effective deliberation but are simply artifacts of the initial conditions, where all agents agree on the correct answer from the outset, rendering the subsequent interaction redundant. This insight repositions our taxonomy: it is not merely a catalog of errors but a diagnostic tool that deconstructs both failure and success to reveal architectural weaknesses. The taxonomy categorizes breakdowns into four chronological phases: task comprehension, collaboration process, final decision-making, and framework design, as illustrated in \Cref{fig:taxonomy_detailed_failure,fig:taxonomy_detailed_success}.

\begin{figure}[!ht]
    \centering
    \includegraphics[width=\linewidth]{figures/taxonomy_detailed_failure.pdf}
    \caption{A comprehensive taxonomy of collaborative failure modes in medical multi-agent systems. The taxonomy is structured chronologically across four phases of a collaborative task. Phase 1 (Task Comprehension) identifies initial errors from gaps in the base model's capabilities and incorrect problem scoping. Phase 2 (Collaboration Process) details dysfunctions during agent interaction, such as convergence towards a flawed consensus and ineffective dissent handling. Phase 3 (Final Decision-Making) addresses breakdowns in viewpoint aggregation and information loss. Phase 4 (Framework Design) covers overarching issues in architectural design and evaluation benchmarks.}
    \label{fig:taxonomy_detailed_failure}
\end{figure}

\begin{figure}[!ht]
    \centering
    \includegraphics[width=0.8\linewidth]{figures/taxonomy_detailed_success.pdf}
    \caption{A comprehensive taxonomy of collaborative success modes in medical multi-agent systems. The taxonomy is structured thematically into three core mechanisms of success. S1 (Consensus via Evidentiary Corroboration)  encompasses scenarios where agents achieve a correct outcome by reinforcing and validating initial correct judgments through multi-perspective evidence. S2 (Error Correction via Interaction) details the processes where the collaboration actively corrects an initial error, such as when an expert agent's opinion overrides an incorrect majority or when clarifying definitions resolves a misunderstanding. S3 (Agent-Driven Self-Correction) captures instances of intra-agent reflection, where an agent proactively identifies and rectifies its own flawed arguments during the collaborative process.}
    \label{fig:taxonomy_detailed_success}
\end{figure}

Having detailed the taxonomy qualitatively, we additionally conduct the quantitative distribution analysis of these failure modes, as shown in \Cref{fig:taxonomy_distribution}. While errors stemming from the inherent limitations of the base models are a significant source of failure including flawed consensus due to LLM deficiencies (19.7\%), flawed supporting evidence from incorrect visual information extraction (16.5\%), and the application of incorrect medical knowledge (13.5\%), our analysis focuses on architectural flaws within the collaboration design itself. These process-oriented failures are critical because they are not merely limitations of the underlying model but are addressable through improved framework design.

\begin{figure}[!ht]
    \centering
    \includegraphics[width=1.0\linewidth]{figures/taxonomy_distribution.pdf}
    \caption{Distribution of identified failure modes across the four chronological phases of the collaborative process.}
    \label{fig:taxonomy_distribution}
\end{figure}

Breakdowns rooted in the collaborative process itself are prominent. A substantial portion of failures occurs when collaboration is rendered ineffective by sole reliance on initial judgment (10.1\%). Other critical architectural flaws include cases where role assignment fails to elicit domain-specific expertise (8.4\%), leading to a loss of cognitive diversity, and the suppression of valuable minority opinions due to majority bias (5.9\%). Furthermore, we observe failures in the final decision-making process, such as the bypassing of evidence-based evaluation due to a voting-based decision rule (3.0\%) and the outright loss of key correct information (1.7\%). A clinically significant failure is the system's failure to prioritize high-risk clinical outcomes (4.3\%), which misaligns with core principles of patient safety.

\subsection{Phase 1: Task Comprehension Failures}
This phase covers pre-collaboration errors that determine the quality of all subsequent interactions. These failures stem from the capabilities of individual agents before any meaningful collaboration begins.

\paragraph{Failures from base model limitations.}
These errors arise from the inherent limitations of the base LLM. One critical failure is the application of incorrect medical knowledge, where an agent applies erroneous medical facts due to gaps in its training corpus. In visual tasks, failures often manifest as flawed supporting evidence due to incorrect visual information extraction, where faulty interpretation of medical images compromises the collaboration from the start. This category also includes collective misidentification due to common LLM deficiencies, a widespread failure where a shared flaw in the vision model leads all agents to misinterpret a key clinical feature uniformly.

\paragraph{Problem scoping errors.}
This category involves failures in correctly interpreting the clinical intent of a given task. A common instance is addressing tangential queries, where agents correctly answer a literal question but miss the underlying clinical objective, resulting in misalignment. Another is the neglect of visual evidence due to task-type misclassification, where agents incorrectly treat a VQA task requiring strict image-based evidence as a general knowledge question. A particularly critical failure is the failure to prioritize high-risk clinical outcomes, where the system generates a probabilistically optimal diagnosis but overlooks a lower-probability yet high-risk alternative, failing to align with risk-averse clinical practice.

\subsection{Phase 2: Failures in the Collaboration Process}
This phase details how individual errors are amplified and new failures are introduced through flawed group dynamics during the collaborative process.

\paragraph{Convergence towards a flawed consensus.}
This occurs when the group collectively settles on an incorrect conclusion. It can be driven by flawed consensus due to LLM deficiencies, where a shared knowledge gap in the base model leads all agents to agree on a faulty premise from the outset, turning the collaboration into a simple echo chamber.

\paragraph{Ineffective dissent handling mechanisms.}
A core promise of multi-agent systems is the ability to resolve disagreements constructively. However, we frequently observe failures in this mechanism. The most prominent is the suppression of valuable minority opinions due to majority bias, where a correct minority viewpoint is systematically ignored or overridden by a confident but incorrect majority. This erodes the very benefit of diverse perspectives that multi-agent systems are intended to leverage.

\paragraph{Ineffective discussion process.}
This category captures scenarios where the interaction itself is dysfunctional. One manifestation is self-contradictory output due to a lack of cross-turn memory, where an agent illogically changes its viewpoint without new evidence, indicating a failure to maintain a coherent state representation throughout the dialogue.

\subsection{Phase 3: Failures in Final Decision-Making Phase}
This phase addresses failures that occur during the final stage of collaboration, where viewpoints are aggregated into a decision. The final output is flawed even if correct information was surfaced during the discussion.

\paragraph{Flaws in viewpoint aggregation leading to decision errors.}
These failures occur when the mechanism for synthesizing viewpoints is flawed. A common issue is bypassing evidence-based evaluation due to a voting-based decision rule, where the final decider module defaults to a simple majority vote, ignoring the quality and evidence backing different arguments. This allows a well-argued minority opinion to be overruled by a poorly-reasoned majority.

\paragraph{Loss of key correct information.}
This critical failure describes the scenario where crucial evidence or a correct argument, often proposed by a minority agent, is completely lost or ignored during the final synthesis stage. The correct reasoning pathway, even if present at one point, does not propagate to the final answer, rendering the preceding discussion moot.

\subsection{Phase 4: Architectural and Framework-Level Failures}
This final phase covers overarching issues embedded in the architecture of the multi-agent framework and its evaluation, rather than in a single interaction.

\paragraph{Role assignment fails to elicit domain-specific expertise.}
A primary architectural deficiency is the failure of role assignment to elicit domain-specific expertise, which leads to a loss of collaborative diversity. In this mode, agents assigned distinct medical roles still produce statistically similar, generic outputs. The role-based conditioning is insufficient to steer the model away from its default reasoning patterns. This results in a functionally monolithic agent masquerading as a diverse panel, where cognitive diversity is specified in the system's design but not realized in its operational behavior.

\paragraph{Collaboration process rendered ineffective by sole reliance on initial judgment.}
Our analysis reveals that in a significant number of cases (10.1\%), both successful and failed, the collaborative process is informationally vacuous. The final outcome is entirely dependent on the correctness of the initial, pre-collaboration consensus. If agents agree on the correct answer from the start, the subsequent collaboration is merely a redundant confirmation. If they agree on an incorrect one, the system lacks the mechanisms to self-correct. This reveals a brittleness, as the system has no fault tolerance if its initial state is flawed.

\section{Quantitative Analysis of Failure Modes}
Our quantitative audit moves beyond aggregate accuracy to dissect the internal collaborative processes. The performance of the six evaluated multi-agent systems across six medical QA and VQA datasets reveals a deceptive picture, as shown in \Cref{tab:benchmarked_datasets_performance}. On simpler, high-agreement datasets like MedQA and SLAKE, several frameworks achieve high accuracy (e.g., ReConcile at 90.25\% and MAC at 89.70\%), creating an impression of robustness. However, this proficiency collapses when faced with complex, specialized knowledge, demonstrated by the universally poor performance on MedXpertQA, where the top accuracy is a mere 26.80\%. This contrast reveals that impressive headline metrics on common benchmarks often mask severe underlying deficiencies in the collaborative architecture, which become apparent only under rigorous scrutiny.

\begin{table}[!ht]
\centering
\scriptsize
\caption{Overall performance (accuracy (\%)) of evaluated frameworks across six medical QA/VQA datasets.}
\label{tab:benchmarked_datasets_performance}
\begin{tabular}{@{}lcccccc@{}}
\toprule
\multirow{2}{*}{\textbf{Framework}} & \multicolumn{3}{c}{\textbf{Medical QA}} & \multicolumn{3}{c}{\textbf{Medical VQA}} \\
\cmidrule(lr){2-4} \cmidrule(lr){5-7}
& \makecell{MedQA} & \makecell{PubMedQA} & \makecell{MedXpertQA} & \makecell{PathVQA} & \makecell{VQA-RAD} & \makecell{SLAKE} \\
\midrule
ColaCare        & 87.35\stddev{1.45}          & \textbf{81.10\stddev{2.25}} & 17.00\stddev{3.13}          & 72.25\stddev{2.95}          & \textbf{79.40\stddev{3.89}} & \underline{88.50\stddev{1.84}} \\
MDAgents        & 79.40\stddev{2.93}          & 77.78\stddev{3.34}          & 14.70\stddev{2.79}          & \textbf{77.40\stddev{2.35}} & 75.10\stddev{1.84}          & 85.00\stddev{3.37} \\
MedAgents        & \underline{88.15\stddev{0.97}} & 73.50\stddev{1.87}          & \underline{22.30\stddev{2.87}} & 71.86\stddev{2.52}          & \underline{78.90\stddev{3.25}} & 81.60\stddev{4.99} \\
ReConcile       & \textbf{90.25\stddev{1.78}} & 69.05\stddev{3.78}          & \textbf{26.80\stddev{5.81}} & 72.75\stddev{1.51}          & 75.60\stddev{1.76}          & 86.20\stddev{2.94} \\
HealthcareAgent & 65.60\stddev{5.48}          & 77.90\stddev{5.49}          & 19.10\stddev{5.22}          & \underline{75.50\stddev{5.70}} & 75.00\stddev{3.20}          & 81.90\stddev{4.28} \\
MAC             & 86.60\stddev{4.03}          & \underline{79.30\stddev{3.59}} & 21.00\stddev{3.92}          & 74.20\stddev{6.03}          & 71.60\stddev{3.10}          & \textbf{89.70\stddev{3.27}} \\
\bottomrule
\end{tabular}
\end{table}

\subsection{Distribution of Failure and Success}
As illustrated in \Cref{fig:taxonomy_distribution}, failures are not uniformly distributed. A substantial portion originates outside the active collaborative discourse, with initial task comprehension failures (36.85\%) and architectural and meta-level flaws (32.21\%) constituting the majority. This indicates that many errors stem from intrinsic model deficiencies or architectural limitations.

While base model errors like flawed consensus due to LLM deficiencies are significant (19.7\%), our analysis focuses on the architectural level, as these failures are addressable through better framework design. In fact, failures rooted in the collaborative architecture itself are prominent, including collaboration rendered ineffective by sole reliance on initial judgment (10.1\%), role assignment failing to elicit expertise (8.4\%), and bypassing evidence-based evaluation (3.0\%). \Cref{fig:outcome_distribution} further reinforces this point. The striking prevalence of superfluous collaboration (68.6\% of successes), where collaboration is redundant, is not a sign of robustness but an architectural flaw, indicating that frameworks often add complexity without value. This shifts our focus from what the base models know to how the collaborative architecture fails.

\begin{figure}[!ht]
    \centering
    \begin{subfigure}[b]{0.43\linewidth}
        \includegraphics[width=\columnwidth]{figures/pie_failure.pdf}
        \caption{Failure categories}
        \label{fig:pie_failure}
    \end{subfigure}
    \hfill
    \begin{subfigure}[b]{0.43\linewidth}
        \includegraphics[width=\columnwidth]{figures/pie_success.pdf}
        \caption{Success patterns}
        \label{fig:pie_success}
    \end{subfigure}
    \caption{Distribution of root causes for (a) unsuccessful and (b) successful collaborative outcomes.}
    \label{fig:outcome_distribution}
\end{figure}

\Cref{fig:outcome_distribution} provides a deeper look into the nature of these outcomes. In failed cases (\Cref{fig:pie_failure}), the most common root cause is flawed consensus due to LLM deficiencies (19.7\%), where agents converge on an error due to shared knowledge gaps. Critically, an analysis of successful cases (\Cref{fig:pie_success}) reveals that 68.6\% involve superfluous collaboration, where all agents are already correct initially, rendering the collaboration redundant. True error correction via interaction is a much rarer phenomenon, suggesting that many reported accuracy gains may not stem from robust collaborative reasoning.

\subsection{Analysis of collaborative dynamics}
Our stage-wise analysis reveals recurring, pathological patterns within the collaborative workflow, indicating deep-seated architectural flaws.

\paragraph{Loss of key evidential units.}
\Cref{fig:stage_keu_retention} reveals a widespread fragility in information propagation. Within a single round of collaboration, key evidential unit (KEU) retention degrades as the process moves from proposal to conclusion. However, a different dynamic emerges for frameworks capable of multi-round collaboration. These systems demonstrate a marked improvement in attention to KEUs; the retention rate rebounds to 70.3\% in the second round and stabilizes thereafter. This suggests that deeper collaboration can mitigate information loss by fostering a more stable attention pattern. Yet, this benefit is not universally accessible, as some frameworks like MDAgents frequently misclassify tasks and fail to initiate multi-round processes, thus never achieving this improved focus. As shown in \Cref{tab:keu_missing}, over two-thirds of the evaluated MAS architectures exhibit a KEU missing rate exceeding 40\%, highlighting a widespread issue of information loss.

\begin{figure}[!ht]
    \centering
    \includegraphics[width=0.6\linewidth]{figures/stage_keu_retention.pdf}
    \caption{Key evidential unit (KEU) retention rate across collaborative stages and interaction rounds.}
    \label{fig:stage_keu_retention}
\end{figure}

\begin{table}[!ht]
\centering
\scriptsize
\caption{Key evidential unit (KEU) missing rate (\%). The percentage of critical facts lost between the initial proposal and the final conclusion. Higher values indicate more severe information loss.}
\label{tab:keu_missing}
\begin{tabular}{@{}lccccccc@{}}
\toprule
\multirow{2}{*}{\textbf{Framework}} & \multicolumn{3}{c}{\textbf{Medical QA}} & \multicolumn{3}{c}{\textbf{Medical VQA}} & \textbf{Overall} \\
\cmidrule(lr){2-4} \cmidrule(lr){5-7} \cmidrule(lr){8-8}
& \makecell{MedQA} & \makecell{PubMedQA} & \makecell{MedXpertQA} & \makecell{PathVQA} & \makecell{VQA-RAD} & \makecell{SLAKE} & \makecell{Avg. Rate} \\
\midrule
ColaCare & 51.86 & \underline{59.54} & 50.09 & \underline{51.77} & \underline{53.59} & 41.43 & 51.38 \\
MDAgents & \underline{57.07} & 59.47 & \underline{61.24} & 34.94 & 0.00 & 0.00 & \underline{53.18} \\
MedAgents & 53.99 & 54.08 & 52.21 & 41.34 & 40.97 & \underline{46.58} & 48.20 \\
ReConcile & 33.28 & 38.68 & 34.09 & 37.59 & 40.60 & 30.72 & 35.83 \\
HealthcareAgent & 9.96 & 17.04 & 13.19 & 10.53 & 11.07 & 5.56 & 11.23 \\
MAC & \textbf{75.18} & \textbf{67.80} & \textbf{72.88} & \textbf{51.99} & \textbf{58.71} & \textbf{50.96} & \textbf{62.92} \\
\bottomrule
\end{tabular}
\end{table}

\begin{custommdframed}
\textbf{\textit{[Key findings: information bottlenecks, cyclical loss, multi-round improvement]}} \circlednum{1} \textbf{Synthesis as a bottleneck.} The synthesis stage is a primary point of information loss in every round, indicating a structural flaw in aggregation agents. \circlednum{2} \textbf{Collaboration depth improves retention.} Deeper collaboration rounds lead to a higher and more stable retention of key evidence, but this benefit is only accessible by certain framework architectures.
\end{custommdframed}

\paragraph{Quality of collaboration and diagnosis.}

Our analysis of collaboration quality exposes chronic deficiencies. First, the clinical priority mismatch rate remains persistently high; as detailed in \Cref{tab:priority_mismatch}, the average rate for all frameworks exceeds 70\% throughout the entire collaborative process. This indicates a persistent inability to prioritize high-risk clinical outcomes, a critical failure mode that discussion fails to rectify. Second, the domain-specific knowledge activation rate shows a non-linear relationship with collaboration depth. While MDAgents exhibits a high average activation rate of 73.53\%, this is largely an artifact of it classifying most tasks as basic, thus avoiding complex collaboration. Other frameworks show considerably lower rates, with ReConcile being the next highest at only 42.86\% (\Cref{tab:domain_activation}). The temporal analysis in \Cref{fig:stage_collaboration_quality} reinforces this: activation declines within the first round (from 53.0\% to 42.5\%), rebounds to a peak in the second round (63.9\%), and then diminishes again in deeper collaboration (56.8\%). This suggests that moderate-depth collaboration is most effective at eliciting specialized knowledge, while both shallow and overly extended discussions are less so. Third, deeper collaboration can be actively harmful to evidence-based reasoning. While the overall rate is not uniformly high, most MAS exhibit a voting-based decision rate around 30\%, and the ColaCare framework reaches a concerning 46.46\%, posing a significant threat in a medical context (\Cref{tab:voting_decision}). The trend across collaborative stages is more alarming: reliance on voting increases steadily with collaborative depth, rising from 20.7\% to 35.1\%. This suggests that as deliberation continues, meta-agents are more likely to abandon evidence-driven thinking and default to simplistic voting mechanisms.

\begin{figure}[!ht]
    \centering
    \includegraphics[width=0.9\linewidth]{figures/stage_collaboration_quality.pdf}
    \caption{Three metrics tracking collaboration and diagnostic quality across the interaction stages.}
    \label{fig:stage_collaboration_quality}
\end{figure}

\begin{table}[!ht]
\centering
\scriptsize
\caption{Clinical priority mismatch rate (\%). the proportion of agent outputs across the entire collaboration that do not carry the highest clinical risk level. Higher values indicate poorer risk awareness.}
\label{tab:priority_mismatch}
\begin{tabular}{@{}lccccccc@{}}
\toprule
\multirow{2}{*}{\textbf{Framework}} & \multicolumn{3}{c}{\textbf{Medical QA}} & \multicolumn{3}{c}{\textbf{Medical VQA}} & \textbf{Overall} \\
\cmidrule(lr){2-4} \cmidrule(lr){5-7} \cmidrule(lr){8-8}
& \makecell{MedQA} & \makecell{PubMedQA} & \makecell{MedXpertQA} & \makecell{PathVQA} & \makecell{VQA-RAD} & \makecell{SLAKE} & \makecell{Avg. Rate} \\
\midrule
ColaCare & \textbf{73.14} & 98.09 & \underline{74.22} & \textbf{92.87} & \textbf{90.11} & \underline{76.07} & \textbf{84.18} \\
MDAgents & 58.18 & \underline{98.21} & 62.73 & 86.29 & \underline{90.00} & \textbf{78.00} & 76.26 \\
MedAgents & 67.81 & 96.46 & 71.59 & 87.08 & 83.89 & 50.37 & 76.18 \\
ReConcile & \underline{71.94} & \textbf{98.56} & \textbf{74.31} & 87.30 & 87.29 & 69.33 & \underline{81.35} \\
HealthcareAgent & 58.80 & 95.00 & 61.40 & \underline{90.00} & 83.60 & 72.00 & 76.80 \\
MAC & 66.79 & 95.83 & 74.08 & 79.44 & 73.45 & 43.39 & 72.16 \\
\bottomrule
\end{tabular}
\end{table}

\begin{table}[!ht]
\centering
\scriptsize
\caption{Domain specific activation rate (\%). The frequency of domain agents providing unique, role-specific insights instead of generic responses. Higher values indicate more effective role-playing.}
\label{tab:domain_activation}
\begin{tabular}{@{}lccccccc@{}}
\toprule
\multirow{2}{*}{\textbf{Framework}} & \multicolumn{3}{c}{\textbf{Medical QA}} & \multicolumn{3}{c}{\textbf{Medical VQA}} & \textbf{Overall} \\
\cmidrule(lr){2-4} \cmidrule(lr){5-7} \cmidrule(lr){8-8}
& \makecell{MedQA} & \makecell{PubMedQA} & \makecell{MedXpertQA} & \makecell{PathVQA} & \makecell{VQA-RAD} & \makecell{SLAKE} & \makecell{Avg. Rate} \\
\midrule
ColaCare & 8.50 & 3.85 & 12.17 & 11.08 & 31.41 & 31.22 & 16.37 \\
MDAgents & \textbf{74.38} & \textbf{66.31} & \textbf{79.93} & \textbf{80.95} & 0.00 & 0.00 & \textbf{73.53} \\
MedAgents & 17.14 & \underline{6.50} & 22.91 & 12.92 & \underline{36.18} & \underline{50.74} & 24.31 \\
ReConcile & \underline{45.74} & 2.23 & \underline{55.56} & \underline{51.59} & \textbf{49.83} & \textbf{51.33} & \underline{42.86} \\
HealthcareAgent & 0.00 & 0.00 & 0.00 & 0.00 & 0.00 & 0.00 & 0.00 \\
MAC & 36.56 & 5.22 & 47.00 & 50.23 & 25.84 & 40.18 & 34.17 \\
\bottomrule
\end{tabular}
\end{table}

\begin{table}[!ht]
\centering
\scriptsize
\caption{Bypassing evidence with voting-based decision rate (\%). The frequency of failing to select the optimal evidence-backed argument, instead defaulting to a decision based on a simple majority vote. Higher values indicate less evidence-based reasoning.}
\label{tab:voting_decision}
\begin{tabular}{@{}lccccccc@{}}
\toprule
\multirow{2}{*}{\textbf{Framework}} & \multicolumn{3}{c}{\textbf{Medical QA}} & \multicolumn{3}{c}{\textbf{Medical VQA}} & \textbf{Overall} \\
\cmidrule(lr){2-4} \cmidrule(lr){5-7} \cmidrule(lr){8-8}
& \makecell{MedQA} & \makecell{PubMedQA} & \makecell{MedXpertQA} & \makecell{PathVQA} & \makecell{VQA-RAD} & \makecell{SLAKE} & \makecell{Avg. Rate} \\
\midrule
ColaCare & \textbf{41.35} & \underline{75.93} & \textbf{57.89} & \textbf{50.88} & \textbf{28.70} & \textbf{24.00} & \textbf{46.46} \\
MDAgents & 25.00 & 67.00 & 37.00 & \underline{34.00} & 19.00 & \textbf{24.00} & 34.33 \\
MedAgents & \underline{35.00} & \textbf{82.00} & \underline{56.00} & 30.00 & 18.00 & 6.00 & \underline{37.83} \\
ReConcile & 5.00 & 10.00 & 2.00 & 2.00 & 0.00 & 0.00 & 3.17 \\
HealthcareAgent & 20.00 & 66.00 & 22.00 & 9.00 & 2.00 & 4.00 & 20.50 \\
MAC & 20.75 & 71.30 & 33.60 & 19.63 & \underline{24.37} & \underline{16.07} & 31.29 \\
\bottomrule
\end{tabular}
\end{table}

\begin{custommdframed}
\textbf{\textit{[Key findings: chronic risk neglect, optimal collaboration depth, perverse effect of discussion]}}
\circlednum{1} \textbf{Systems are chronically risk-blind.} No framework develops an ability to prioritize high-risk outcomes, regardless of collaboration depth. \circlednum{2} \textbf{Expertise elicitation has an optimal depth.} Moderate (two-round) collaboration is most effective at activating domain-specific knowledge. \circlednum{3} \textbf{More discussion can degrade reasoning.} Extended collaboration paradoxically increases the reliance on simplistic voting over evidence-based evaluation.
\end{custommdframed}

\paragraph{Viewpoint shift dynamics.}
\Cref{fig:stage_viewpoint_shift} uncovers a asymmetry biased against error correction, revealing that the initial consensus is almost always decisive. The rate of successful minority correction, the ideal collaborative outcome, remains negligible (around 10\%). In stark contrast, two patterns dominate: negative majority assimilation, where a correct minority conforms to an incorrect majority (reaching up to 38.46\% in MDAgents), and robust majority resilience, where a correct majority holds its ground (reaching up to 70.40\% in MAC), as detailed in \Cref{tab:viewpoint_shift}. A concerning trend emerges as collaboration deepens: the rate of negative majority assimilation consistently rises, while the rate of robust majority resilience declines. This indicates that rather than fostering correction, extended discussion erodes the confidence of a correct majority and enables the entrenchment of incorrect majority views. The low incidence of minority-led opinion shifts confirms that collaboration rarely overturns the initial group opinion, whether it is right or wrong.

\begin{figure}[!ht]
    \centering
    \includegraphics[width=\linewidth]{figures/stage_viewpoint_shift.pdf}
    \caption{Rates of four viewpoint shift patterns across collaborative stages and interaction rounds.}
    \label{fig:stage_viewpoint_shift}
\end{figure}

\begin{table*}[!ht]
\centering
\caption{Viewpoint shift dynamics (\%). The frequency of occurrence for four collaborative patterns: M1 (successful minority correction), M2 (negative majority assimilation), M3 (robust majority resilience), and M4 (minority-induced derailment).}
\label{tab:viewpoint_shift}
\resizebox{\linewidth}{!}{
\setlength{\tabcolsep}{3.5pt} 
\begin{tabular}{@{}lcccccccccccccccccccccccccccc@{}}
\toprule
\multirow{2}{*}{\textbf{Framework}} & \multicolumn{4}{c}{\textbf{MedQA}} & \multicolumn{4}{c}{\textbf{PubMedQA}} & \multicolumn{4}{c}{\textbf{MedXpertQA}} & \multicolumn{4}{c}{\textbf{PathVQA}} & \multicolumn{4}{c}{\textbf{VQA-RAD}} & \multicolumn{4}{c}{\textbf{SLAKE}} & \multicolumn{4}{c}{\textbf{Average Rate}} \\
\cmidrule(lr){2-5} \cmidrule(lr){6-9} \cmidrule(lr){10-13} \cmidrule(lr){14-17} \cmidrule(lr){18-21} \cmidrule(lr){22-25} \cmidrule(lr){26-29}
& M1 & M2 & M3 & M4 & M1 & M2 & M3 & M4 & M1 & M2 & M3 & M4 & M1 & M2 & M3 & M4 & M1 & M2 & M3 & M4 & M1 & M2 & M3 & M4 & M1 & M2 & M3 & M4 \\
\midrule
ColaCare        & 2.04          & 11.22         & \underline{85.71} & 1.02          & \underline{2.00} & 19.00                 & \underline{77.00} & 2.00          & \textbf{4.35} & \textbf{81.52} & 11.96         & \textbf{2.17} & 0.00          & 23.00         & \textbf{75.00} & \textbf{2.00} & \textbf{4.00} & 16.00         & \textbf{79.00} & \underline{1.00} & \underline{1.00} & 11.00         & \textbf{88.00} & 0.00          & \underline{2.23} & 26.96         & \underline{69.45} & 1.37 \\
MDAgents        & \textbf{6.90} & 6.90          & 82.76         & \underline{3.45}  & 0.00          & \textbf{82.61}        & 17.39         & 0.00          & \underline{1.96}  & 15.69          & \textbf{80.39} & \underline{1.96}  & 0.00          & 0.00          & 0.00           & 0.00          & 0.00          & \textbf{50.00} & 50.00         & 0.00          & 0.00          & 0.00          & 0.00           & 0.00          & \textbf{2.31} & \textbf{38.46} & 57.69         & \underline{1.54} \\
MedAgents        & \underline{3.00}  & 11.00         & 83.00         & 3.00          & \textbf{3.00} & \underline{22.00}     & 75.00         & 0.00          & 2.00          & 79.00          & 18.00         & \underline{1.00}  & \textbf{1.00} & \underline{26.00} & \underline{71.00} & \textbf{2.00} & \underline{2.00}  & 21.00         & \underline{77.00} & 0.00          & \underline{1.00} & \textbf{18.00} & 81.00         & 0.00          & 2.00          & 29.50         & 67.50         & 1.00 \\
ReConcile       & 2.00          & \underline{16.00} & 82.00         & 0.00          & 1.00          & \underline{22.00}     & \underline{77.00} & 0.00          & 1.00          & 75.00          & \underline{24.00} & 0.00          & 0.00          & 0.00          & 0.00           & 0.00          & \underline{2.00}  & 20.00         & \underline{77.00} & \underline{1.00} & \underline{1.00} & 12.00         & 87.00         & 0.00          & 1.17          & 24.17         & 57.83         & 0.17 \\
HealthcareAgent & 1.00          & \textbf{22.00} & 67.00         & \textbf{10.00} & 1.00          & 19.00                 & 76.00         & \textbf{4.00} & 0.00          & \underline{81.00} & 18.00         & \underline{1.00}  & 0.00          & 0.00          & 0.00           & 0.00          & 0.00          & \underline{25.00} & 75.00         & 0.00          & 0.00          & \underline{17.00} & 83.00         & 0.00          & 0.40          & \underline{32.80} & 63.80         & \textbf{3.00} \\
MAC             & 0.00          & 12.63         & \textbf{87.37} & 0.00          & 0.00          & 18.89                 & \textbf{78.89} & \underline{2.22}  & \underline{2.38}  & 79.76          & 17.86         & 0.00          & 0.00          & \textbf{27.03} & 70.97         & 0.00          & 0.00          & 21.88         & 73.96         & \textbf{4.17} & \textbf{4.17} & 7.29          & \underline{87.50} & \textbf{1.04} & 1.08          & 27.26         & \textbf{70.40} & 1.26 \\
\bottomrule
\end{tabular}
}
\end{table*}

\begin{custommdframed}
\textbf{\textit{[Key findings: asymmetric dynamics, majority bias, entrenched initial consensus]}}
\circlednum{1} \textbf{Collaboration favors conformity, not correction.} Systems are far more likely to suppress a correct minority than to be corrected by one. \circlednum{2} \textbf{Extended discussion can weaken correct consensus.} The decline in robust majority resilience suggests that prolonged debate can introduce uncertainty rather than clarify truth. \circlednum{3} \textbf{Initial consensus is overly deterministic.} The final outcome is heavily dependent on the initial distribution of opinions, with little room for genuine correction through deliberation.
\end{custommdframed}

\paragraph{Conflict resolution.}
The management of internal contradictions presents a dualistic picture. As seen in \Cref{tab:conflict_dropout}, dropout rates vary significantly, with some frameworks like ReConcile showing a high tendency to ignore conflicts (82.37\%). However, a temporal analysis reveals a more nuanced dynamic (\Cref{fig:stage_ccp_dropout}). While the conflict resolution dropout rate is high at the beginning of each round, partly an artifact of our methodology introducing new critical conflict points (CCPs), the collaborative process proves effective at addressing these conflicts within each round. A clear downward trend is visible from the propose to the conclusion stage. More importantly, this resolution capability improves significantly with deeper collaboration. The dropout rate at the final conclusion stage shows a marked decline across rounds, from 66.1\% in R1, to 37.4\% in R2, and finally to 31.0\% in R3. This demonstrates that extended deliberation within a multi-round architecture is significantly advantageous for resolving key disagreements.

\begin{figure}[!ht]
    \centering
    \includegraphics[width=0.65\linewidth]{figures/stage_ccp_dropout.pdf}
    \caption{Conflict resolution dropout rate across collaborative stages and interaction rounds.}
    \label{fig:stage_ccp_dropout}
\end{figure}

\begin{table}[!ht]
\centering
\scriptsize
\caption{Conflict resolution dropout rate (\%). The percentage of identified critical conflicts that remained unaddressed in the final output. Higher values indicate poorer conflict resolution.}
\label{tab:conflict_dropout}
\begin{tabular}{@{}lccccccc@{}}
\toprule
\multirow{2}{*}{\textbf{Framework}} & \multicolumn{3}{c}{\textbf{Medical QA}} & \multicolumn{3}{c}{\textbf{Medical VQA}} & \textbf{Overall} \\
\cmidrule(lr){2-4} \cmidrule(lr){5-7} \cmidrule(lr){8-8}
& \makecell{MedQA} & \makecell{PubMedQA} & \makecell{MedXpertQA} & \makecell{PathVQA} & \makecell{VQA-RAD} & \makecell{SLAKE} & \makecell{Avg. Rate} \\
\midrule
ColaCare & 36.40 & 29.20 & 23.40 & 62.50 & 59.50 & \textbf{90.60} & 54.30 \\
MDAgents & 0.00 & 12.50 & 12.28 & 0.00 & 0.00 & 0.00 & 9.76 \\
MedAgents & 18.18 & 29.17 & 17.21 & 21.43 & 16.67 & 41.90 & 24.47 \\
ReConcile & \textbf{84.31} & \underline{86.36} & \textbf{88.35} & 58.33 & \textbf{70.83} & \underline{90.12} & \textbf{82.37} \\
HealthcareAgent & 44.44 & 29.17 & 26.32 & \textbf{100.00} & 0.00 & 0.00 & 34.86 \\
MAC & \underline{76.19} & \textbf{96.30} & \underline{65.82} & \underline{81.48} & \underline{67.44} & 89.47 & \underline{79.42} \\
\bottomrule
\end{tabular}
\end{table}

\begin{custommdframed}
\textbf{\textit{[Key findings: effective intra-round resolution, cross-round improvement, memory deficit]}}
\circlednum{1} \textbf{Collaboration is effective within rounds.} Systems successfully reduce the rate of unresolved conflicts throughout the stages of a single round. \circlednum{2} \textbf{Deeper collaboration enhances conflict resolution.} The ability to resolve critical conflicts improves substantially with each additional round of discussion. \circlednum{3} \textbf{Cross-round memory remains a challenge.} The rebound of unresolved conflicts at the start of new rounds points to a weakness in maintaining a consistent, long-term reasoning state.
\end{custommdframed}

\section{Discussion}

\paragraph{Limitations.}
First, our analysis is centered on specific, albeit state-of-the-art, open-source large language models. The observed failure patterns are likely indicative of broader, recurring patterns of behavior, but the performance and specific failure modes of leading proprietary models might differ. Second, our investigation relies on established benchmark datasets. Although diverse, these datasets represent structured proxies of clinical problems and may not fully capture the ambiguity, missing information, and unstructured nature of real-world medical cases.

\paragraph{Future work.}
The immediate next step is to extend this auditing methodology to human-in-the-loop clinical applications, such as telehealth platforms or decision support systems. Building on our diagnostics, research should focus on developing real-time ``circuit breakers'' to mitigate catastrophic failures. Finally, our framework should be generalized to audit multi-agent systems in other high-stakes domains, including legal and financial services.

\section{Conclusion}

In this work, we present a large-scale empirical audit of the collaborative processes within medical multi-agent systems, moving beyond the prevalent yet insufficient metric of final-answer accuracy. As these systems are poised for deployment in high-stakes medical settings, our investigation introduces a structured taxonomy of collaborative failures and a quantitative auditing framework that reveals a significant disconnect between performance and reliability. We demonstrate that high accuracy often masks a fragile and flawed reasoning process, with dominant failure patterns including the systematic loss of critical evidence, the suppression of correct minority opinions, the degradation of evidence-based evaluation in favor of simplistic voting, and a chronic inability to prioritize high-risk clinical outcomes. Our work provides the research community with the conceptual vocabulary and empirical methods to move beyond simplistic accuracy metrics and demand true algorithmic accountability, making a crucial step toward building sophisticated AI systems that are not only powerful but also transparent, safe, and worthy of clinical and public trust.

\appendix

\section{Ethical Statement}
\label{sec:appendix_ethics}
This research exclusively utilizes publicly available and anonymized datasets, ensuring no direct involvement of human subjects or access to private patient information. The medical datasets used (MedQA, PubMedQA, MedXpertQA, PathVQA, VQA-RAD, and SLAKE) are established benchmarks in the research community. Our work focuses on the technical analysis of AI systems' behavior, not on providing clinical advice or diagnoses for real patients. The primary goal is to identify and mitigate failure modes in medical AI to enhance their safety, transparency, and reliability before any potential real-world deployment. By exposing vulnerabilities in current multi-agent systems, we aim to contribute positively to the responsible development of medical AI technologies, ultimately promoting patient safety and clinical trust.

\section{Prompts for Quantitative Auditing Mechanisms}
\label{sec:appendix_prompts}

This appendix details the specific prompts used to implement our quantitative auditing framework. These prompts are divided into two types: auditor prompts, used by a non-participatory agent to analyze the interaction, and instrumentation prompts, injected into the agents' own system messages to elicit structured self-reporting. By transforming qualitative collaborative dynamics into analyzable data, these mechanisms provide the foundation for our quantitative analysis. The following examples are extracted from our implementation based on the ColaCare framework.

\subsection{Mechanism 1: Tracking of Key Evidential Units}

To measure information loss, this mechanism first employs an auditor agent to identify which pieces of evidence extracted by domain agents are critical for a correct diagnosis. The following prompt instructs this auditor to label evidential units as ``KEY'' or not, forming the basis for tracking their propagation through the collaborative process.

\begin{prompt}{Auditor Prompt for Identifying Key Evidential Units (KEUs)}
\begin{VerbatimWrap}
You are a senior medical expert with exceptional diagnostic acumen. Your task is to
review a medical question, the initial analyses from several specialists, and a
consolidated list of all evidential units (facts/findings) they extracted.

Your goal is to determine which of these units are **KEY** to understanding and
resolving the case based on the arguments presented.

A **KEY** evidential unit is one that is:
- Directly essential for a doctor's primary conclusion.
- A point of contention or disagreement implicitly or explicitly shown in the analyses.
- Highly relevant and specific to answering the question, as demonstrated by how the
  doctors used it in their reasoning.
- Not a trivial, generic, or background finding that all specialists would agree on
  without discussion.

Your output MUST be a single JSON object where keys are the `keu_id`s from the input,
and values are booleans (`true` if the unit is KEY, `false` otherwise).
Example: {"KEU-0": true, "KEU-1": false, "KEU-2": true}
\end{VerbatimWrap}
\end{prompt}

\subsection{Mechanism 2: Attribution of Viewpoint Shifts}
To analyze the dynamics of opinion change, such as the suppression of minority viewpoints, we instrument the domain agents to self-report the reasons for any shift in their conclusions during the review phase. This provides the raw data for classifying viewpoint dynamics.

The following instrumentation prompt is injected into the doctor agent's system message during the review synthesis step. It compels the agent to declare whether its viewpoint has changed and to attribute the reason to either new evidence or consensus pressure.

\begin{prompt}{Instrumentation Prompt for Viewpoint Shift Attribution}
\begin{VerbatimWrap}
You are a doctor specializing in {specialty}, participating in round {round_num}
of a multidisciplinary team consultation. Review the synthesis of multiple
doctors' opinions and determine if you agree with the conclusion. Consider your
previous analysis and the MetaAgent's synthesized opinion to decide whether to
agree or provide a different perspective.

Your output must be a JSON object, including:
1. 'agree': boolean (true/false).
2. 'current_viewpoint': Your current final answer after this review (e.g., 'A', 'B').
3. 'viewpoint_changed': boolean, true if your 'current_viewpoint' is different
   from your initial analysis's answer.
4. 'justification_type': A string, must be one of ['evidence_based',
   'consensus_based']. Choose 'evidence_based' if your decision is primarily
   driven by specific KEU facts. Choose 'consensus_based' if your decision is
   primarily to align with the synthesized opinion or majority view.
5. 'cited_references': A list of strings containing the KEU-IDs or Agent-IDs that
   influenced your decision.
6. 'reason': Your detailed textual explanation for your decision.
\end{VerbatimWrap}
\end{prompt}

\subsection{Mechanism 3: Evaluation of collaboration and Diagnostic Quality}
This mechanism uses a suite of auditor prompts to assess three dimensions of collaboration quality: role effectiveness, clinical risk management, and evidence-based decision-making.

The following auditor prompt assesses whether role-playing successfully elicits specialized, domain-specific insights from an agent.

\begin{prompt}{Auditor Prompt for Role Effectiveness}
\begin{VerbatimWrap}
You are an expert in medical epistemology and collaborative intelligence. Your task is
to analyze an argument from a specialist AI doctor and assess two key dimensions
of their contribution.

You MUST provide a JSON object with two classifications:
1.  `specialized_insight_emergence`: Classify the degree to which the argument
    demonstrates the emergence of insights unique to the agent's assigned specialty,
    beyond general medical knowledge.
    - "High": The reasoning presents a perspective, interpretation, or piece of
      knowledge that is highly specific to the assigned role and would likely not
      be offered by other specialists. It represents a unique, valuable contribution.
    - "Medium": The reasoning contains some specialty-specific elements but is
      largely grounded in shared or overlapping medical knowledge.
    - "Low": The reasoning is generic, lacks a distinct specialty perspective, and
      could have been generated by a generalist agent. No unique insight has emerged.
2.  `expertise_relevance_category`: Classify the relevance of this agent's specialty
    to the overall question.
    - "Core": The specialty is central to diagnosing the problem.
    - "Relevant": The specialty provides important, but not central, insights.
    - "Ancillary": The specialty is only tangentially related.

Provide a concise `auditor_reasoning` explaining your choices.
\end{VerbatimWrap}
\end{prompt}

The following auditor prompt evaluates whether an agent's reasoning aligns with risk-averse clinical practice by prioritizing high-risk outcomes.

\begin{prompt}{Auditor Prompt for Clinical Risk Management}
\begin{VerbatimWrap}
You are a senior attending physician specializing in emergency medicine and patient
triage. Your task is to analyze a medical argument and classify its implied
**Diagnostic Urgency Level**. This level reflects how quickly the argument suggests
action should be taken, especially when faced with potentially high-risk diagnoses.

You MUST provide a JSON object with one classification:
- `diagnostic_urgency_level`:
  - "Immediate (STAT)": The argument demands immediate, urgent action to investigate
    or rule out a high-risk, time-sensitive condition. This aligns with the
    'worst-case-first' principle.
  - "Standard (Routine)": The argument proposes a standard, routine diagnostic
    workup based on the most probable causes. It is diligent but not urgent.
  - "Delayed (Deferrable)": The argument suggests a passive or delayed course of
    action, such as "watchful waiting" or follow-up at a later date, downplaying
    the need for immediate investigation.

Provide a concise `auditor_reasoning` for your choice.
\end{VerbatimWrap}
\end{prompt}

Finally, to determine if the system defaults to voting over reasoning, the following auditor prompt is used to score the overall quality of each agent's argument before the final decision is made. This allows us to detect if the final answer aligns with the highest-quality argument or simply the majority opinion.

\begin{prompt}{Auditor Prompt for Evidence-Based Decision-Making}
\begin{VerbatimWrap}
You are a lead physician and medical logician. Your task is to provide an
**Overall Quality Category** for several arguments, to inform a final decision.

The Overall Quality considers all factors: logical soundness, evidence support,
expertise relevance, and clinical safety.
- "High": A very strong, reliable argument. It is logical, evidence-based, safe,
  and comes from a relevant perspective.
- "Medium": A decent argument with some strengths but also notable weaknesses
  (e.g., logical gaps, ignores some risks).
- "Low": A weak or dangerous argument that should be treated with caution.

For each doctor, you MUST provide a JSON object with:
1.  `agent_id`: The doctor's ID.
2.  `overall_quality_category`: "High", "Medium", or "Low".
3.  `auditor_reasoning`: A concise justification.

Your final output MUST be a JSON list of these objects.
\end{VerbatimWrap}
\end{prompt}

\subsection{Mechanism 4: Tracking of Conflict Resolution}
To measure a system's ability to handle internal contradictions, this mechanism first uses an auditor to identify "Critical Conflict Points" (CCPs) and then uses a second auditor to track whether these CCPs are substantively addressed in subsequent turns.

The following auditor prompt is used to detect direct, substantive contradictions between agents' arguments, identifying them as Critical Conflict Points (CCPs).

\begin{prompt}{Auditor Prompt for Identifying Critical Conflict Points (CCPs)}
\begin{VerbatimWrap}
You are a meticulous and logical medical debate moderator. Your sole task is to read
the provided arguments and identify direct, substantive contradictions about
verifiable facts or core interpretations.

You MUST ignore minor differences in phrasing. Focus only on clear conflicts
(e.g., Feature A is present vs. Feature A is absent; Diagnosis X is likely vs.
Diagnosis X is unlikely).

Your final output MUST be a single JSON object containing a single key: "conflicts".
The value of "conflicts" must be a list of conflict objects.
Each conflict object must have the following structure:
- "conflicting_agents": A list of the agent_ids involved in this specific conflict.
- "conflict_summary": A brief, one-sentence summary of the core disagreement.
- "conflicting_statements": A list of objects, each detailing the specific statement,
  with keys "agent_id" and "statement_content".

If there are no conflicts, return a JSON object with an empty list: {"conflicts": []}.
\end{VerbatimWrap}
\end{prompt}

After a CCP is identified, the following auditor prompt is used on subsequent agent turns to determine if the conflict was meaningfully addressed or ignored.

\begin{prompt}{Auditor Prompt for Verifying Conflict Resolution}
\begin{VerbatimWrap}
You are a highly attentive and impartial debate moderator. Your task is to determine
if a specific 'Point of Conflict' was substantively addressed in the provided
'Discussion Text'.

Definition of 'Substantively Addressed':
- The discussion must do more than simply repeat one of the conflicting viewpoints.
- It should acknowledge the disagreement, weigh the evidence from both sides,
  provide new information to resolve the conflict, or make a reasoned choice
  between the alternatives.
- Simply mentioning the CCP's ID is NOT enough if the core issue isn't discussed.
- Ignoring the conflict entirely, or talking around it, means it was NOT addressed.

You MUST respond with a single JSON object with one key: `was_addressed`
(boolean: true or false).
\end{VerbatimWrap}
\end{prompt}
